\section{初步应用}\label{sec:ag-2-3}
% 原PDF p32，§2.3起.
\begin{numpara}[有理映射的延拓 1]\label{item:ag-2-3-1}
设 $X$ 是概形，$\mathcal L$ 是可逆层，$s_0,\ldots,s_n$ 是若干截面，不要求它们生成 $\mathcal L$. 记 $X_{s_i}$ 为截面 $s_i$ 不为零的开集；在此开集上，$s_i$ 给出平凡化 $s_i:\mathcal O_{X_{s_i}}\xrightarrow\sim\mathcal L|_{X_{s_i}}$. 于是得到 $n$ 个截面 $s_j/s_i\in\Gamma(X_{s_i},\mathcal O_X)$，从而得到态射 $X_{s_i}\longrightarrow\mathbb A^n=D(T_i)\subset\mathbb P^n$. 这些态射粘合为 $U:=\bigcup_iX_{s_i}\longrightarrow\mathbb P^n$.
不能将此态射延拓到 $X$，但在爆破之后可以延拓. 事实上，设 $\mathcal F\subset\mathcal L$ 是由 $s_i$ 生成的 $\mathcal O_X$-子模. 它不再可逆，但在理想 $\mathcal F\otimes\mathcal L^{\otimes-1}\subset\mathcal O_X$ 的爆破 $\widetilde X\xrightarrow\pi X$ 上变得可逆. 因而，$\pi^*\mathcal L$ 的截面 $\pi^*s_i$ 生成一个可逆子模，定义态射 $\widetilde X\longrightarrow\mathbb P^n$. 由构造，$\pi$ 诱导同构 $\pi^{-1}(U)\xrightarrow\sim U$；将 $U$ 与 $\pi^{-1}(U)$ 等同后，如此构造的态射便延拓了原来的态射.
\end{numpara}
\begin{numpara}[有理映射的延拓 2]\label{item:ag-2-3-2}
设 $\mathcal I$ 是 $X$ 上由截面 $f_0,\ldots,f_n$ 生成的理想层. 设 $U$ 为 $\mathcal I$ 定义的闭集的补集. 这是 $f_i$ 不同时为零的点所成的集合. 因此，它们定义态射 $U\longrightarrow\mathbb P^n$. 同样，此态射不能延拓到 $X$，但可以延拓到 $X$ 在 $\mathcal I$ 处的爆破 $\widetilde X_{\mathcal I}$. 从某种意义上说，这次爆破是为使此态射能够延拓而须对 $X$ 作出的最小改造. 更确切地说，有如下命题.
\end{numpara}
% 原PDF p33.
\begin{proposition*}
沿用上述记号，态射 $\widetilde X_{\mathcal I}\longrightarrow\mathbb P^n\times X$ 是闭浸入，并将 $\widetilde X_{\mathcal I}$ 等同于 $U\longrightarrow\mathbb P^n$ 的图像的概形论闭包.
\end{proposition*}
\begin{proof}
所考虑的态射是闭浸入，因为它由 $\mathcal O_X[T_0,\ldots,T_n]\longrightarrow\operatorname{Rees}^{\mathcal I}(\mathcal O_X)$ 诱导，而后者由构造是满射. 此态射在 $U$ 上的限制是 $U\longrightarrow\mathbb P^n$ 的图像. 又因 $U$ 是一个除子的补集，它在 $\widetilde X_{\mathcal I}$ 中概形论稠.
\end{proof}
由证明可知，甚至可以将 $U$ 换成其中一个概形论稠的开集. 也可以从 $\mathcal I$ 的亚纯截面，即 $\mathcal I\otimes\mathcal M_X$ 的截面出发. 下面给出一个应用.
\begin{example}\label{item:ag-2-3-3}
设 $S$ 是离散赋值谱，其环为 $R$，一致化元为 $\varpi$，泛点为 $\eta$，闭点为 $s$（可设想 $R=\mathbb Z_p$）. 记 $0_{\varpi}$ 为 $\mathbb P_S^1$ 的特殊纤维 $\mathbb P_s^1$ 中的点 $[0:1]$. 于是，$0_{\pi}$ 的理想由亚纯截面 $x/y$ 与 $\varpi$ ``生成''. 因而，$\mathbb P_S^1$ 在 $0_{\pi}$ 处的爆破可以实现为映射 $[x:y]\in\mathbb P_\eta^1\mapsto[x:\varpi y]\in\mathbb P_\eta^1$ 的图像在 $\mathbb P_S^1\times\mathbb P_S^1$ 中的概形论闭包 $\widetilde{\mathbb P^1}$，即局部闭子集
\begin{align*}
\{([x:y],[x:\varpi y]),\ \text{其中 }[x:y]\in\mathbb P_\eta^1\}
\end{align*}
的概形论闭包. 更确切地说，爆破 $\pi_{0_{\varpi}}$ 可等同于复合 $\widetilde{\mathbb P^1}\hookrightarrow\mathbb P_S^1\times\mathbb P_S^1\xrightarrow{\pi_1}\mathbb P_S^1$. 现在可看出，$\widetilde{\mathbb P^1}$ 也是 $[x,y]\in\mathbb P_\eta^1\mapsto([\varpi x,y],[x,y])\in\mathbb P_\eta^1\times\mathbb P_\eta^1$ 的像的闭包. 因而，复合 $\widetilde{\mathbb P^1}\hookrightarrow\mathbb P_S^1\times\mathbb P_S^1\xrightarrow{\pi_2}\mathbb P_S^1$ 可等同于 $\mathbb P_s^1$ 中的点 $\infty_{\varpi}=[1,0]$ 的爆破. 于是有图
\begin{align*}
\begin{tikzcd}[ampersand replacement=\&,column sep=large]
\& \widetilde{\mathbb P^1}\arrow[dl,"\pi_{0_{\varpi}}"']\arrow[dr,"\pi_{\infty_{\varpi}}"] \& \\
\mathbb P_S^1\arrow[rr,dashed,"{[x:y]\mapsto[x:\varpi y]}"'] \& \& \mathbb P_S^1.
\end{tikzcd}
\end{align*}
可以给出 $\widetilde{\mathbb P^1}$ 的模解释. 事实上，态射
\begin{align*}
\begin{tikzcd}[ampersand replacement=\&,column sep=large]
\mathcal O_{\mathbb P^1}(-1)^2\arrow[r,"{\left(\begin{smallmatrix}1&0\\0&\varpi\end{smallmatrix}\right)}"] \& \mathcal O_{\mathbb P^1}(-1)^2\arrow[r,two heads] \& \mathcal O_{\mathbb P^1}
\end{tikzcd}
\end{align*}
（由万有满态射 $\mathcal O_{\mathbb P^1}^2\twoheadrightarrow\mathcal O_{\mathbb P^1}(1)$ 导出）的像恰好是定义点 $0_{\varpi}$ 的理想. 利用概形 $\widetilde{\mathbb P^1}$ 是此理想的爆破这一事实，可知它分类如下图：
\begin{align*}
\begin{tikzcd}[ampersand replacement=\&]
\mathcal O^2\arrow[r]\arrow[d,two heads] \& \mathcal O^2\arrow[d,two heads] \\
\mathcal L'\arrow[r] \& \mathcal L
\end{tikzcd}
\end{align*}
其中 $\mathcal L$ 与 $\mathcal L'$ 可逆，上方态射由矩阵 $\left(\begin{smallmatrix}1&0\\0&\varpi\end{smallmatrix}\right)$ 给出. 到 $\mathbb P_S^1$ 的爆破态射 $\pi_{0_{\varpi}}$ 和 $\pi_{\infty_{\varpi}}$ 分别对应于忘掉此图的其中一列.\unskip\nobreak\hfill\ensuremath{\diamond}
\renewcommand{\envenddiamond}{}% marker on the last text line
\end{example}
% 原PDF p34.
\emph{本节余下内容未在课堂上讲授.}
\begin{numpara}[群作用的延拓]\label{item:ag-2-3-4}
给定整概形 $X$，以及有限（抽象）群 $G$ 在 $X$ 的非空开集 $U$ 上的作用. 当 $X$ 仿射时，容易看出，此作用可延拓到 $X$ 的正规化上，而正规化（在实际情形中，即 $X$ 优良时）在 $X$ 上有限. 但当 $X$ 不仿射时，情形就复杂了.
回到前面的例子，考虑离散赋值谱 $S$ 上的 $\mathbb P^1$. 考察群 $\mathbb Z/2\mathbb Z$ 在泛纤维 $\mathbb P_\eta^1$ 上由 $\varphi:[x,y]\mapsto[\varpi y,x]$ 给出的作用（即矩阵 $\left(\begin{smallmatrix}0&\varpi\\1&0\end{smallmatrix}\right)$）. 态射 $\varphi$ 不能延拓到整个 $\mathbb P_S^1$，但可延拓到特殊纤维 $\mathbb P_s^1$ 中闭点 $0_{\varpi}$ 的补集. 因而必须爆破此点. 事实上，这就足够了，如下图所示，其中用到了分解
\begin{align*}
\begin{gathered}
\begin{pmatrix}0&\varpi\\1&0\end{pmatrix}=\begin{pmatrix}0&1\\1&0\end{pmatrix}\begin{pmatrix}1&0\\0&\varpi\end{pmatrix}.\\[4pt]
\begin{tikzcd}[ampersand replacement=\&,column sep=large]
\& \widetilde{\mathbb P^1}\arrow[dl,"\pi_{0_{\varpi}}"']\arrow[dr,"\pi_{\infty_{\varpi}}"]\arrow[rr,dotted] \& \& \widetilde{\mathbb P^1}\arrow[dr,"\pi_{0_{\varpi}}"] \& \\
\mathbb P_S^1\arrow[rr,dashed,"{[x:y]\mapsto[x:\varpi y]}"'] \& \& \mathbb P_S^1\arrow[rr,"{[x:y]\mapsto[y:x]}"'] \& \& \mathbb P_S^1
\end{tikzcd}
\end{gathered}
\end{align*}
点线箭头的存在性来自爆破的万有性质. 由泛纤维的概形论稠性，其平方自动为恒等态射.
\end{numpara}
\begin{numpara}[不可约分支的分离]\label{item:ag-2-3-5}
设 $Y$ 是 $X$ 的闭子概形，它是 $X$ 的两个闭子概形 $Z$ 与 $Z'$ 的并，且二者互不包含. 设 $\widetilde X$ 是沿 $Z\cap Z'$ 的爆破. 则严格变换 $\widetilde Y$ 是严格变换 $\widetilde Z$ 与 $\widetilde Z'$ 的不交并.
事实上，记 $V$（分别为 $V'$）为 $Z$（分别为 $Z'$）在 $X$ 中的开补集，记 $U=V\cup V'$ 为爆破 $\pi$ 在其上为同构的开集. 再记 $\widetilde V$ 为 $\widetilde X$ 中使例外除子 $E$ 局部由定义 $Z$ 的理想的一个正则截面生成的开集. 同样定义 $\widetilde V'$. 于是有 $\widetilde X=\widetilde V\cup\widetilde V'$ 及 $\widetilde V\setminus E=\pi^{-1}(V)$. $Z$ 与 $Z'$ 互不包含的假设蕴含 $\widetilde V$ 与 $\widetilde V'$ 也互不包含. 已知 $\widetilde Z$ 可等同于 $\pi^{-1}(U\cap Z)$ 在 $\widetilde X$ 中的概形论闭包. 因而，它的支集包含在 $\pi^{-1}(U\setminus V)$ 的拓扑闭包中，从而包含在 $\widetilde X\setminus\widetilde V'$ 中. 同理，$\widetilde Z'$ 的支集包含在 $\widetilde X\setminus\widetilde V$ 中. 而
\begin{align*}
(\widetilde X\setminus\widetilde V)\cap(\widetilde X\setminus\widetilde V')=\varnothing.
\end{align*}
\end{numpara}
